%% kapa: exam=2 round=35 year=2024 period=1 subject=appraisal_practice subject_ko=감정평가실무 seq=3 number=3 points=20
%% source: Q-Net 감정평가사 기출문제 (cst00309) artlSeq=5213532 file=raw/1008/5213532/2024년 제35회 감정평가사 2차시험 문제지.pdf
%% conversion: kordoc 4.24.51 -> tools/gen2.py (auto); review=proofread-sonnet; flags=none
\begin{longquestion}{【문제 3】}
도시개발사업이 시행중인 구역 내의 토지 기호(1) 및 기호(2)에 대하여 해당 ‘토지대장’과 ‘환지확정처분조서’에 근거하여 < 자료 1 >과 같이 정리하였다. < 자료 2 >의 가격자료를 참고하여 다음 물음에 답하시오. (20점) \par
\medskip\noindent 물음 1) 청산금이 정산(교부 또는 납부)된 상태를 전제로 기호(1) 및 기호(2) 토지의 과도 또는 부족면적을 판정하여 사정면적을 확정하고, 청산금의 정산 상태를 고려한 현재의 가격을 각각 추정하여 면적의 차이를 분석하시오. (10점) \par
\medskip\noindent 물음 2) 청산금이 미정산된 상태를 전제로 기호(1) 및 기호(2) 토지의 과도 또는 부족면적을 판정하여 사정면적을 확정하고, 청산금의 미정산 상태를 고려한 현재의 가격을 각각 추정하여 면적의 차이를 분석하시오. (10점) \par
\medskip\noindent < 자료 1 > 토지의 상황 \par
{\small\centering
\resizebox{\linewidth}{!}{%
\begin{tabular}{|c|c|c|c|c|c|c|c|c|c|c|}
\hline
\multirow{2}{*}{구분} & \multicolumn{4}{c|}{종전토지 (토지대장기준)} & \multicolumn{4}{c|}{종후의 환지예정지} & \multirow{2}{*}{\makecell{과도\\면적\\(m\textsuperscript{2})}} & \multirow{2}{*}{\makecell{부족\\면적\\(m\textsuperscript{2})}} \\
\cline{2-9}
 & 소재지 & 지번 & 지목 & \makecell{면적\\(m\textsuperscript{2})} & BL & LT & \makecell{권리\\면적\\(m\textsuperscript{2})} & \makecell{환지\\면적\\(m\textsuperscript{2})} &  &  \\
\hline
기호(1) & \makecell{A시\\B동} & \makecell{10\\번지} & 전 & 600 & 35 & 13 & 420 & 460 & - & - \\
\hline
기호(2) & \makecell{A시\\B동} & \makecell{20\\번지} & 답 & 1,200 & 35 & 14 & 840 & 800 & - & - \\
\hline
\end{tabular}}\par}

\medskip\noindent < 자료 2 > 가격자료 \par
1. A시 B동에서 기호(1) 및 기호(2)와 유사한 지목 ‘전’과 ‘답’ 등의 농경지 가격은 100,000원/m\textsuperscript{2} 수준에 형성되어 있다. \par
2. A시 B동에서 도시개발사업이 기 완료된 기호(1) 및 기호(2)와 유사한 환지된 토지의 가격은 200,000원/m\textsuperscript{2} 수준에 형성되어 있다. \par
3. 본 도시개발사업과 관련한 교부 또는 납부(징수) 청산금의 단위면적 당 가격은 도시개발사업 완료 후 가격과 유사한 200,000원/m\textsuperscript{2} 수준이다. \par
4. 현재의 시점은 사업의 종료가 임박한 추상적 시점으로 시점차이에 대한 보정의 필요성은 없는 것으로 조사되고, 기호(1)과 기호(2)의 소유자는 상이하여 상계처리대상이 아니다. \par
\end{longquestion}
